% R3 — cross-disease transfer and the adaptation curve (synced to P0 round-17 state, 2026-09-09)
\subsection{Cross-disease transfer is unresolvable in one direction and rank-indeterminate in the other}
\label{ssec:r3_cross_disease}

We then tested whether representations learned in one disease carry to the
other, training on COVID-19 and testing on sepsis and vice versa under the
primary severe-as-survived label definition, with patient-level scoring and
patient-clustered bootstrap intervals (Section~\ref{ssec:methods_transfer}).
Forward transfer (COVID$\rightarrow$sepsis; 91 sepsis patients, 17 fatal) is
uninformative by construction: Section~\ref{ssec:r2_within_sepsis} establishes
that mortality is not learnable in this pool even with in-domain training, so
a transfer measurement into it re-measures that null rather than transfer
ability. Point estimates spanned 0.644 (\model{Geneformer}) to 0.684
(\model{UCE}), with bootstrap confidence intervals overlapping throughout
(Table~\ref{tab:transfer} and Figure~\ref{fig:transfer}). Reverse transfer
(sepsis$\rightarrow$COVID; 96 sepsis training patients, 240 COVID-19 test
patients, 25 fatal) tests the direction where signal exists and gave higher
point estimates throughout: \model{STATE} 0.721 (95\% CI [0.621, 0.809]) and
\model{Geneformer} 0.712 (95\% CI [0.612, 0.814]) with logistic regression
(0.711 and 0.636 with random forest), followed by \model{UCE} (0.693/0.607),
\model{scGPT} (0.642/0.622), and \model{scPRINT-2} (0.619/0.517) (logistic
regression/random forest), with all bootstrap intervals overlapping. No
reverse-direction ranking can be defended either: the severe-excluded
sensitivity analysis reverses the top two models (\model{Geneformer} 0.775,
\model{STATE} 0.765), and adding only five training patients moves
\model{UCE} by $+0.086$ AUC between the 91- and 96-patient training pools
(Table~\ref{tab:transfer}). Transfer estimates are therefore
not identifiable at this scale, which is itself the strongest evidence the
experiment provides. The forward/reverse asymmetry is a property of the two
target tasks, one learnable and one not, not of transfer direction.

To ask whether transfer can be recovered with target-disease supervision, we
measured the adaptation curve: COVID-trained classifiers received increasing
numbers of sepsis training cells and were scored on a held-out 32-patient
sepsis test set, a one-third stratified holdout of the 96-patient labeled
sepsis pool (Section~\ref{ssec:methods_transfer}). Even with all
target-disease training data, all five foundation models and the PCA-200
baseline remained between 0.40 and 0.62 patient-level AUC (zero-shot values
lower still), with no foundation model separating from the baseline at any
training size (Figure~\ref{fig:adaptation_curve}). We report this as
unresolvable rather than null: at roughly 3 fatal events in 32 test patients,
the 95\% interval around any patient-level AUC spans about 0.7, so the
observed range is consistent with a single value plus noise, and no
rescue-or-not conclusion is identifiable at this test-set size. Two
comparability notes accompany this result. First, cross-disease
transfer on the 5-patient GSE167363 sepsis test set (2 fatal) produced
substantially higher point estimates (up to 1.00 patient-level AUC); given
the test-set size, we report those values
only as a labeled small-sample reference (Supplementary Note~\ref{sn:gse167363_testset}). Second, two additional
models evaluated under a different test split (TranscriptFormer and STACK) are
excluded from this comparison for the same reason and are noted in the same
supplementary record.

Together, the transfer experiments provide no evidence for a transferable
inflammatory-mortality representation on this underrepresented cell type:
forward transfer is uninformative, reverse transfer is rank-indeterminate,
and the adaptation curve is below the resolution of its test set.
